\documentclass[11pt]{article}
\newcommand{\tablelang}{uk}
\input{../../../paper/tex/preamble}
\newcolumntype{L}{>{\raggedright\arraybackslash}X}
\newcolumntype{R}{>{\raggedleft\arraybackslash}X}
\newcolumntype{P}[1]{>{\raggedright\arraybackslash}p{#1}}
\crefname{section}{розділ}{розділи}
\Crefname{section}{Розділ}{Розділи}
\crefname{appendix}{додаток}{додатки}
\Crefname{appendix}{Додаток}{Додатки}
\setlength{\emergencystretch}{2em}
\hypersetup{pdftitle={Коли оцінка fitness вводить в оману: емпіричне дослідження еволюції програм під керуванням LLM},
  pdfauthor={Oleksandr Pylypenko},
  pdfkeywords={еволюція програм під керуванням LLM, OpenEvolve, пошук програм, валідність оцінювача, надійність бенчмарку, Gemma}}

\begin{document}
\selectlanguage{ukrainian}
\title{Коли оцінка fitness вводить в оману: емпіричне дослідження еволюції програм під керуванням LLM}
\author{Oleksandr Pylypenko\\[2pt]\small Незалежний дослідник у галузі штучного інтелекту та машинного навчання\\
\small Аспірант Державного торговельно-економічного університету, Київ, Україна\\
\small У співпраці з кафедрою інженерії програмного забезпечення та кібербезпеки\\[2pt]
\small ORCID: \href{https://orcid.org/0009-0004-9610-9242}{0009-0004-9610-9242}}
\date{19 листопада 2026 року}
\maketitle

\begin{center}\small\textit{Чернетка препринту. Рукопис не подавався на рецензування та не розміщувався в архіві препринтів.}\end{center}

\begin{abstract}
Еволюційний пошук програм під керуванням великих мовних моделей (LLM) використовує оцінювач і для спрямування пошуку, і для вибору кандидата. Якщо оцінювач не враховує вимогу задачі, додатковий пошук може посилити цю невідповідність, а не усунути її. У статті описано два попередньо зареєстровані дескриптивні дослідження з OpenEvolve 0.4.0 на задачах пакування кіл і наближеного пошуку найближчих сусідів. У дослідженні масштабу на Colab L4 порівнювали конфігурації Gemma 4 E4B, 12B, 26B A4B і 31B, еволюційний пошук, best-of-N, збільшені бюджети та частково виконаний режим із хмарною допомогою. Для еволюційного пошуку медіани незалежно перевіреної оцінки пакування кіл становили 1.63, 1.63, 2.17 та 2.25 відповідно. Загальний тест між групами моделей був статистично значущим (Kruskal--Wallis $p=0.015$), однак парні порівняння не пройшли корекцію Голма й не доводять еквівалентності. Після серійного перерахунку kNN хмарно підтриманий пошук мав вищу оцінку у восьми з дев'яти парних запусків ($p=0.0078$, розвідувальний аналіз). Водночас 28\% його кандидатів, оцінених робочим оцінювачем вище за початкову програму, виявилися невалідними або не швидшими за неї під час незалежної перевірки. У MoE-прогоні на 1{,}000 викликів валідні кандидати досягли 2.43, але робочий оцінювач вибрав невалідну програму у двох із трьох запусків. Результати характеризують невідповідність оцінювача у двох обмежених задачах; вони не встановлюють загальних законів масштабування чи причинного ефекту хмарної допомоги.
\end{abstract}

\noindent\textbf{Ключові слова:} еволюційний пошук програм під керуванням LLM; OpenEvolve; пошук програм; валідність оцінювача; надійність бенчмарку; незалежна перевірка.

\medskip
\noindent\textbf{Незалежність дослідження.} Експерименти виконав автор, використовуючи власне обладнання та облікові записи сервісів. Робота не координувалася з Google, супроводжувачами OpenEvolve чи авторами цитованих досліджень і не схвалена ними. Автор повідомляє про академічний статус і співпрацю, зазначені вище.

% =====================================================================
\section{Вступ}\label{sec:intro}

Еволюційний пошук програм під керуванням LLM поєднує механізм пропонування кандидатів з оцінювачем. Модель пропонує код; оцінка за результатом виконання спрямовує відбір; успішні кандидати стають контекстом для наступних пропозицій. Системи FunSearch і AlphaEvolve демонструють потенціал такого підходу для задач із виконуваною цільовою функцією та цінними алгоритмічними результатами~\citep{romera2024funsearch,novikov2025alphaevolve}. Відкриті реалізації поширили цей дизайн на ширше коло моделей і задач~\citep{assumpcao2025codeevolve,sharma2026openevolve}.

Оцінювач є частиною ефективної специфікації пошуку. Кандидат може отримати високу оцінку, водночас порушивши вимогу, якої немає у функції оцінювання. Під час багаторазового пошуку зростає ймовірність такого випадку: обраний для використання кандидат має найбільшу оцінку, а не обов'язково є найкоректнішим. Ця проблема пов'язана з проблемою тестового оракула у програмній інженерії та результатами про те, що звичайні набори тестів можуть пропускати функціональні помилки згенерованих програм~\citep{barr2015oracle,liu2023evalplus}.

Ми досліджуємо це питання у двох зареєстрованих дескриптивних дослідженнях OpenEvolve на двох невеликих задачах. У дослідженні Colab змінювали конфігурації Gemma 4, порівнювали еволюційний пошук (L) із best-of-N семплюванням (B), використовували режим еволюції з хмарною допомогою (LC) і збільшували бюджет пошуку. У пілоті на ноутбуці порівнювали L і B з моделлю 12B. Обидві задачі мають робочий оцінювач, який спрямовує пошук, та незалежну перевірку після завершення. Це дає змогу відокремити оцінку пошукового процесу від перевіреного результату.

Наш внесок полягає в емпіричному аудиті результатів пошуку та узгодженості оцінювача за цих умов. Ми не пропонуємо нового еволюційного алгоритму. Попередні праці вже оцінювали OpenEvolve та споріднені системи на тому самому бенчмарку пакування 26 кіл; ця робота зосереджується на хибних покращеннях, незалежній перевірці, спостереженнях за різних масштабів моделей, впливі бюджету та наслідках шуму вимірювання часу~\citep{assumpcao2025codeevolve}. Ставимо такі дослідницькі питання:

\begin{itemize}
\item \textbf{RQ1: масштаб моделі та режим пошуку.} Як змінюються перевірені оцінки пакування кіл для чотирьох конфігурацій Gemma 4 та між L і B за однакового бюджету у 100 викликів?
\item \textbf{RQ2: бюджет пошуку.} Чи покращує довший пошук найкращого валідного кандидата та оцінку кандидата, обраного робочим оцінювачем?
\item \textbf{RQ3: хмарна допомога.} Як LC і L порівнюються у парних MoE-прогонах на пакуванні кіл і kNN?
\item \textbf{RQ4: узгодженість оцінювача.} Як часто кандидати, оцінені робочим оцінювачем вище за початкову програму, не проходять незалежну перевірку або не покращують перевірену оцінку початкової програми?
\end{itemize}

Попередні реєстрації передбачали дескриптивне звітування без тестів статистичної значущості. Усі інференційні порівняння в цьому рукописі додано після кампаній і прямо позначено як розвідувальні. Ми подаємо їх як підсумок цих запусків, а не як підтверджувальні тести чи оцінки загальних причинних ефектів.

% =====================================================================
\section{Пов'язані дослідження}\label{sec:related}

\paragraph{Еволюційний пошук під керуванням LLM.} FunSearch поєднує генератор-кандидатів на основі LLM із виконанням програм та еволюційним циклом пошуку математичних і алгоритмічних розв'язків~\citep{romera2024funsearch}. AlphaEvolve поширює підхід на наукове й промислове відкриття алгоритмів, застосовуючи один або кілька оцінювачів для керування ітеративними змінами коду~\citep{novikov2025alphaevolve}. Відкриті фреймворки дають змогу запускати такі методи з різними моделями. CodeEvolve містить контрольовані порівняння з OpenEvolve за узгоджених конфігурацій і бюджетів, зокрема на CirclePackingSquare з $n=26$~\citep{assumpcao2025codeevolve}. У цій роботі основна увага приділена якості розв'язків та компонентам алгоритму. Ми ж звітуємо, як часто ранжування робочого fitness-оцінювача узгоджується з незалежною перевіркою в іншому режимі моделей і виконання.

\paragraph{Сила оцінювання та тестові оракули.} Проблема оракула полягає у визначенні того, як тест встановлює коректність результату, коли очікуване значення недоступне безпосередньо~\citep{barr2015oracle}. EvalPlus показав, що посилення тестів у бенчмарках генерації коду може виявляти помилки у програмах, раніше зарахованих як правильні~\citep{liu2023evalplus}. Наше середовище має спільні риси, але відрізняється: система пошуку має власну функцію оцінювання, а незалежна перевірка застосовується після відбору кандидата. Ми вимірюємо розбіжність і на рівні кандидатів, і на рівні запусків.

\paragraph{Ефективність за умови коректності.} ECCO оцінює методи генерації та редагування коду для підвищення ефективності програм, водночас контролюючи функціональну коректність; автори також звертають увагу на надійність часових вимірювань на різному обладнанні~\citep{waghjale2024ecco}. Наша задача kNN має подібну проблему вимірювання, але використовує ітеративний еволюційний цикл і чутливу до швидкості робочу оцінку. Серійне повторне вимірювання в цій статті коригує конкретне джерело шуму, але не замінює ширший бенчмарк ефективності.

\paragraph{Реалістичність задач.} Бенчмарки ремонту репозиторіїв, зокрема SWE-bench, оцінюють, чи виправляє згенерований моделлю патч реальну проблему GitHub у наявній кодовій базі~\citep{jimenez2024swebench}. У цьому дослідженні використовуються невеликі обмежені задачі з явними виконуваними оцінювачами. Такий вузький дизайн полегшує незалежну перевірку кандидатів, але зовнішню валідність ще потрібно дослідити на реальних репозиторіях.

% =====================================================================
\section{Дизайн і метрики}\label{sec:design}

\subsection{Режими пошуку}

Усі кампанії використовували OpenEvolve 0.4.0 із зареєстрованими початковими програмами та конфігураціями~\citep{sharma2026openevolve}. У режимі \textbf{L} Gemma створює мутації в еволюційній популяції та архіві. У режимі \textbf{B} та сама модель незалежно семплює кандидатів від початкової програми без еволюційної історії; робочий оцінювач ранжує кандидатів. У режимі \textbf{LC} зберігається еволюційний цикл, приблизно 10\% викликів мутації спрямовується до Gemini 4 Argon, а решта виконується Gemma. Бюджет визначає кількість зареєстрованих викликів LLM, а не спільну одиницю обчислень, затримки чи енергії.

\subsection{Задачі та робочі оцінювачі}

\textbf{Пакування кіл.} Задача полягає в розміщенні 26 кіл в одиничному квадраті без перетинів і максимізації суми радіусів. Робочий оцінювач використовує нормовану суму радіусів (поділену на 2.635) і надає нуль кандидатам, які відхиляє. Незалежна перевірка вимагає скінченних додатних радіусів, розміщення всередині квадрата, неперетину за строгішою похибкою, однакового результату у двох запусках sandbox і дотримання статичних обмежень. У статті подаємо сиру суму радіусів за повторною перевіркою; відома верхня межа оцінки становить близько 2.635. Початкова програма сама є невалідною, оскільки два радіуси дорівнюють нулю, тому її перевірена оцінка дорівнює нулю.

\textbf{Наближений пошук $k$ найближчих сусідів.} Робочий оцінювач поєднує повноту пошуку та відносну швидкість порівняно з точним пошуком, але не перевіряє, чи повернуті квадрати відстаней відповідають індексам повернутих сусідів. Незалежна перевірка використовує нові тестові входи, перевіряє запити-копії та правильність відстаней, а швидкість і повноту оцінює на окремих відкладених наборах. Отже, кандидат може отримати високу робочу оцінку завдяки скороченню точності, але не пройти вимогу задачі щодо відстаней.

\subsection{Незалежні кінцеві показники}

Для запуску $r$ позначимо через $s_w(c)$ оцінку робочого оцінювача для кандидата $c$, а через $s_v(c)$ -- оцінку незалежної перевірки, якщо кандидат проходить вимоги валідності. Обраний кандидат $c_r^*=\arg\max_c s_w(c)$. Основний кінцевий показник -- оцінка повторної перевірки; якщо кандидат не проходить перевірку валідності, запуску зараховується перевірена оцінка початкової програми. \emph{Хибним покращенням} вважається оцінка кандидата, для якого $s_w(c)>s_w(c_0)$, але він невалідний за повторної перевірки або має $s_v(c)\leq s_v(c_0)$, де $c_0$ -- початкова програма. Об'єднаний знаменник хибних покращень рахує оцінювання кандидатів, включно з повторними оцінюваннями того самого хешу коду, а не лише унікальні програми.

\begin{figure}[tb]
\centering\includegraphics[width=\linewidth]{fig1_pipeline_uk.pdf}
\caption{Конвеєр пошуку й оцінювання. Режими L, LC і B генерують набори кандидатів. Робочий оцінювач керує відбором, а незалежна перевірка визначає опублікований перевірений результат. Відкладена перевірка недоступна генератору кандидатів та механізму відбору за робочою оцінкою.}
\label{fig:pipeline}
\end{figure}

\subsection{Кампанії та відхилення протоколу}

Дослідження на Colab було зареєстроване до початку і передбачало 179 запусків. Завершено 167; відсутні 12 були прогонами MoE LC, які не запустили або зупинили через обмеження квоти хмарної моделі. Стандартна шкала за бюджетом у 100 викликів мала по десять seeds на режим для E4B, 12B і 26B A4B на пакуванні кіл і kNN, та по п'ять seeds на режим для 31B лише на пакуванні кіл. Додатково виконано три MoE L-прогони по 1{,}000 викликів для пакування кіл, три MoE L-прогони по 500 викликів для kNN і три MoE B-прогони по 1{,}000 викликів для пакування кіл; довгий режим B є розвідувальним. Для LC завершено по дев'ять парних запусків на кожну задачу.

Окремий пілот на ноутбуці передбачав 18 запусків і завершив 12: L і B на кожній задачі з трьома seeds на режим. Режим LC закрили після вичерпання квоти; незавершений запуск виключено. У ноутбуці використовувалася відеокарта RTX 4060 Laptop із 8 GB VRAM та 31 GB RAM. На Colab використовувався один L4 із 24 GB VRAM та 12 vCPU. Інші відхилення та деталі протоколу описані в двох попередніх реєстраціях і журналах відхилень~\citep{pylypenko2026artifacts}.

Запуски 12B на Colab і ноутбуці не утворюють контрольованого порівняння обладнання. Вони використовували різні збірки моделі, версії Ollama та середовища виконання кандидатів. У режимі Colab 12B B збірка багаторазово генерувала незастосовний блок пошуку й заміни; у запусках пакування кіл кандидатів не було створено. Це відоме обмеження слід враховувати при тлумаченні контрасту L--B для 12B.

\subsection{Корекція вимірювання часу kNN}

Перший перерахунок kNN запускав кілька Docker sandbox одночасно. За такого навантаження та сама початкова програма отримувала оцінки від 0.02 до 1.0, оскільки функція оцінювання враховує швидкість за настінним часом. Тому ми повторно виміряли час для початкової програми та всіх валідних кандидатів вище за неї, виконуючи їх послідовно, по одному sandbox, на незавантаженому ноутбуці. Для хешу початкової програми є вісім серійних вимірювань: п'ять призначених калібрувальних та ще три з запусків, у яких той самий код вибрали як програму для доставки. Основний кінцевий показник об'єднує всі вісім (медіана 0.2235; діапазон 0.1606--0.2438); лише п'ять калібрувальних мають медіану 0.238 (діапазон 0.1919--0.2438). У пілотному L показник хибних покращень змінюється з 36/42 у чутливісному аналізі лише за калібрувальними вимірюваннями до 33/42 за об'єднаною медіаною. Для LC kNN значення залишається 60/213, а LC перевищує L у 8/9 пар за обох визначень.

Серійний перерахунок охопив 214 унікальних валідних хешів кандидатів вище за початкову оцінку з Colab та ще 14 хешів пілоту. Близько 1{,}900 кандидатів нижче за робочу оцінку початкової програми не вимірювали повторно. Тому кінцевий показник найкращої перевіреної програми серед усіх кандидатів не подаємо. Для 71 хешу з повторними серійними вимірюваннями медіанний розмах оцінки становив 0.018; у 16 він перевищував 0.05, максимальний був 0.490. Перерахунок покращує порівнюваність виміряних кандидатів, але залишкова варіативність зберігається біля верхньої межі оцінки.

\subsection{Статистичний аналіз}

Обидві реєстрації визначали лише дескриптивний аналіз. Наведені після завершення досліджень двосторонні порівняння позначено як розвідувальні. Результати пакування кіл у чотирьох групах моделей для режиму L узагальнено тестом Kruskal--Wallis, а парні порівняння -- тестами Mann--Whitney з корекцією Голма. Порівняння L--B також наведені окремо для задач і моделей, із корекцією Голма на дев'ять наявних порівнянь. Для LC--L використано парні seeds, парний знаковий тест, тест Вілкоксона зі знаковими рангами та процентильні бутстреп-інтервали на 10{,}000 повторних вибірках парних різниць. Ці аналізи не встановлюють еквівалентності, причинних ефектів чи загальної ефективності поза цими задачами.

% =====================================================================
\section{Результати}\label{sec:results}

\subsection{Оцінки пакування кіл для різних конфігурацій моделей}

У таблиці~\ref{tab:scale} наведено медіани перевіреної оцінки для 100 викликів. Для режиму L медіани становили 1.63 для E4B і 12B, 2.17 для 26B A4B MoE і 2.25 для 31B. Загальна відмінність між групами моделей була виявлена розвідувальним тестом Kruskal--Wallis ($p=0.0155$). Жодне з трьох парних порівнянь MoE не пройшло корекцію Голма: скориговане $p=0.063$ для E4B, $p=0.173$ для 12B і $p=0.426$ для 31B. Зокрема, результат MoE--31B не доводить еквівалентності.

\begin{table}[tb]
\centering\small
\caption{Перевірені результати пакування кіл у шкалі моделей із бюджетом у 100 викликів. Оцінка -- сира сума радіусів; валідність визначається незалежною перевіркою. Значення $p$ для L--B є розвідувальними та скоригованими за Голмом у групі з дев'яти порівнянь задач і моделей.}
\label{tab:scale}
\begin{tabularx}{\linewidth}{@{}lrrrrr@{}}
\toprule
Модель & Медіана L & Медіана B & $n$ на режим & L--B $p$ & Голм $p$ \\
\midrule
Gemma 4 E4B & 1.628 & 0.851 & 10 & 0.491 & 1.000 \\
Gemma 4 12B & 1.633 & 0.000 & 10 & 0.0022 & 0.0199 \\
Gemma 4 26B A4B MoE & 2.170 & 1.911 & 10 & 0.0373 & 0.299 \\
Gemma 4 31B & 2.249 & 2.400 & 5 & 0.151 & 1.000 \\
\bottomrule
\end{tabularx}
\end{table}

Контраст 12B на пакуванні кіл не є свідченням загальної переваги L над B: збірка Colab B не застосовувала diff, тому всі десять запусків залишилися на невалідній початковій програмі. Для 31B медіана B у цій вибірці з п'яти seeds перевищила L. Інші медіани мають дескриптивний характер; ці результати не підтримують універсальну перевагу L над B. Один 100-викличний запуск на Colab коштував близько 0.9 compute units для MoE і 3.9 для 31B. Модель 31B обслуговувалась одним запуском за раз, тоді як інші моделі виконували по два паралельних запуски. Тому це приблизне платформне порівняння витрат, а не вимірювання енергії чи ефективності за однакового throughput. E4B використала дещо менше compute units на одну медіанну одиницю оцінки (близько 0.40 проти 0.43 для MoE), але мала нижчу абсолютну медіану.

\begin{figure}[tb]
\centering\includegraphics[width=0.92\linewidth]{fig2_circle_scale_uk.pdf}
\caption{Результати окремих запусків пакування кіл для режимів L і B у шкалі моделей із бюджетом 100 викликів. Точки є незалежно перевіреними результатами; темні горизонтальні риски показують медіани. Оцінки 12B B дорівнюють нулю, оскільки збірка Colab не застосувала жодного кандидатного diff.}
\label{fig:scale}
\end{figure}

\subsection{Довший пошук і валідність обраної програми}

Три MoE-прогони L на 1{,}000 викликів знайшли валідні кандидати з оцінками 2.432, 2.204 і 2.258. Однак робочий оцінювач вибрав невалідну програму у двох запусках, тому основні результати обраної програми становили 2.432, 0 та 0. У розвідувальному контролі B на 1{,}000 викликів обрані кандидати мали оцінки 1.936, 2.072 і 1.924 та були валідними. Частка хибних покращень у MoE L long становила 607/1{,}064 кандидатних оцінювань (57\%); у B long -- 358/2{,}089 (17\%). Це невелике порівняння показує, що великий пул кандидатів може містити сильні валідні розв'язки, а робочий оцінювач усе одно може вибрати невалідний. Воно не доводить, що довгий еволюційний пошук загалом гірший за best-of-N.

У MoE L-прогонах kNN на 500 викликів робочий оцінювач не поставив жодного кандидата вище за початкову програму, тому всі три запуски залишили початкову програму. Зареєстрований кінцевий показник найкращої перевіреної програми серед усіх кандидатів неможливо відновити із серійного перерахунку, бо не всі кандидати нижче за робочу оцінку початкової програми повторно вимірювалися.

\subsection{Запуски MoE з хмарною допомогою}

У таблиці~\ref{tab:lc} порівнюються дев'ять парних запусків LC--L для кожної задачі. У пакуванні кіл оцінка LC була вищою у семи парах, однаковою в одній і нижчою в одній. Медіани становили 2.630 для LC і 2.192 для L; розвідувальний тест Вілкоксона не був значущим ($p=0.148$). Для kNN після серійного перерахунку LC мав вищу оцінку у восьми з дев'яти пар. Медіана становила 1.000 проти 0.223 для L, а медіанна парна різниця -- 0.776 (95\% бутстреп-інтервал 0.192--0.777). Двосторонній розвідувальний знаковий тест без корекції дав $p=0.039$, а тест Вілкоксона -- $p=0.0078$. Ці порівняння не були попередньо зареєстровані й не визначають загального причинного ефекту хмарної допомоги. У 18 завершених LC-запусках було 198 викликів Argon і 4.64 мільйона токенів.

\begin{table}[tb]
\centering\small
\caption{Парні результати MoE LC--L. Кожне порівняння має дев'ять seeds. Усі інференційні тести post hoc і без корекції.}
\label{tab:lc}
\begin{tabularx}{\linewidth}{@{}lrrrrL@{}}
\toprule
Задача & Медіана L & Медіана LC & \shortstack{LC вище / \\ порівну / нижче} & Wilcoxon $p$ & Примітка \\
\midrule
Пакування кіл & 2.192 & 2.630 & 7 / 1 / 1 & 0.148 & У малій вибірці немає надійного напряму \\
kNN, серійний час & 0.223 & 1.000 & 8 / 0 / 1 & 0.0078 & Розвідувально \\
\bottomrule
\end{tabularx}
\end{table}

\begin{figure}[tb]
\centering\includegraphics[width=0.95\linewidth]{fig3_lc_pairs_uk.pdf}
\caption{Парні результати MoE для режимів L лише з Gemma та LC з хмарною допомогою. Лінія з'єднує один і той самий seed. Наведені значення $p$ отримано post hoc непарним тестом Вілкоксона зі знаковими рангами без корекції.}
\label{fig:lc}
\end{figure}

\subsection{Пілот на ноутбуці}

На ноутбуці L видав валідне пакування кіл у всіх трьох запусках (медіана 1.773; діапазон 1.596--2.120). B видав невалідну програму у двох із трьох запусків, а медіана основної оцінки становила нуль (діапазон 0--1.793). За трьох seeds це лише дескриптивний результат; розвідувальний тест Mann--Whitney дав $p=0.376$. У всіх шести пілотних запусках kNN кандидат із найвищою робочою оцінкою не пройшов незалежну перевірку, тому L і B повернулися до оцінки початкової програми 0.223. Програми повертали сусідів із високою повнотою, але використовували арифметику float32, що порушувало зазначену в специфікації точність справжніх квадратів відстаней. Робочий оцінювач цю вимогу не перевіряв.

Результати 12B на ноутбуці й Colab не дають змоги відокремити ефект обладнання. Крім різних середовищ виконання, відрізнялися збірки моделі. Збій Colab 12B B, описаний вище, є конкретним прикладом взаємодії збірки моделі з циклом пошуку на основі diff.

\subsection{Хибні покращення за незалежної перевірки}

У таблиці~\ref{tab:false} наведено, що в кожній умові з кандидатами вище за початкову оцінку були хибні покращення, хоча їх частка залежала від задачі, режиму й бюджету. У MoE LC kNN 60 із 213 оцінювань кандидатів вище за початкову програму були хибними (28\%): чотири кандидати не пройшли перевірку валідності, а 56 були валідними, але не швидшими за початкову програму з об'єднаною серійною медіаною. У пілотному L kNN хибними були 33/42 оцінювань (79\%), а в B -- усі п'ять із п'яти кандидатів вище за початкову оцінку. Оскільки для пілотного B $n=5$, його частка має високу невизначеність.

\begin{table}[tb]
\centering\small
\caption{Хибні покращення серед оцінювань кандидатів, які робочий оцінювач ранжував вище за початкову програму. Записи: хибні / оцінювання вище за початкову (частка). Тире означає, що жоден кандидат не отримав оцінки вище за початкову.}
\label{tab:false}
\begin{tabularx}{\linewidth}{@{}lrr@{}}
\toprule
Умова дослідження & Пакування кіл & kNN \\
\midrule
E4B L & 239 / 505 (47\%) & 1 / 2 (50\%) \\
E4B B & 213 / 408 (52\%) & 0 / 1 (0\%) \\
12B L, Colab & 103 / 423 (24\%) & 6 / 9 (67\%) \\
12B B, Colab & --- & 1 / 5 (20\%) \\
26B MoE L & 129 / 421 (31\%) & 1 / 1 (100\%) \\
26B MoE B & 117 / 710 (17\%) & --- \\
26B MoE LC & 140 / 486 (29\%) & 60 / 213 (28\%) \\
MoE L, довгий бюджет & 607 / 1{,}064 (57\%) & --- \\
MoE B, довгий бюджет & 358 / 2{,}089 (17\%) & не запускали \\
12B ноутбук L & 21 / 140 (15\%) & 33 / 42 (79\%) \\
12B ноутбук B & 22 / 86 (26\%) & 5 / 5 (100\%) \\
\bottomrule
\end{tabularx}
\end{table}

Хибні покращення в пакуванні кіл часто виникали через різницю між робочою оцінкою та строгішими критеріями валідності, зокрема через кола з нульовим радіусом. У kNN робочий оцінювач не перевіряв коректність повернутих відстаней і винагороджував швидкі наближені програми. Частка хибних покращень зросла до 57\% у довгих MoE L-прогонах пакування кіл. Це узгоджується з тим, що триваліший пошук відкриває недоліки оцінювача, але порівняння є розвідувальним; довгі запуски відрізнялися бюджетом і розподілом кандидатів.

% =====================================================================
\section{Обговорення}\label{sec:discussion}

\paragraph{Fitness є операційною ціллю, а не доказом коректності.} В обох задачах кандидати могли отримати від робочого оцінювача вищу оцінку за початкову програму і не пройти незалежну перевірку. Випадок kNN особливо показовий: повноти пошуку та швидкості недостатньо, якщо контракт API вимагає також точних квадратів відстаней. Перш ніж нараховувати бонус за швидкість, оцінювач має перевірити коректність і обмеження задачі. Надійний дизайн може використовувати коректність як гейт до вимірювання продуктивності, хоча це дослідження не порівнювало альтернативні дизайни оцінювача експериментально.

\paragraph{Бюджет пошуку може посилити дефекти оцінювача.} Довгі MoE L-прогони знайшли сильні валідні пакування кіл, але робочий оцінювач обрав невалідних кандидатів у двох із трьох запусків. Отже, більша кількість викликів збільшила пул найкращих валідних кандидатів, але не гарантувала кращого вибраного результату. Майбутні дослідження мають окремо звітувати про найкращу валідну програму та перевірений результат кандидата, обраного системою, а також аудіювати саме той код, який пошук передає далі.

\paragraph{Масштаб моделі та обчислювальні витрати.} Спостережувані медіани режиму L зростають між меншими конфігураціями та MoE/31B, але парні порівняння не проходять зазначену корекцію. Вибірка для 31B має п'ять запусків; тест MoE--31B не встановлює ні відмінності, ні еквівалентності. Порівняння compute units також відображає різну паралельність обслуговування. Ці дані мотивують масштабніші парні дослідження з однаковим способом обслуговування, хешами збірок моделей та обліком енергії або грошових витрат.

\paragraph{Хмарна допомога.} У MoE LC перевірена оцінка kNN перевищила L у восьми з дев'яти пар після корекції часу. Post hoc характер аналізу, невелика кількість пар, неповне виконання запланованих LC-прогонів і додаткові хмарні виклики обмежують силу цього результату. У пакуванні кіл результати LC також були спрямовані на його користь, але лишилися непевними. Наступне дослідження має попередньо зареєструвати порівняння, повторити його на незалежних задачах і моделях та обчислити хмарну вартість у спільній одиниці.

\paragraph{Внесок статті.} Попередні праці показали, що еволюційний пошук під керуванням LLM може знаходити високоякісні програми, зокрема в родині задач пакування кіл~\citep{romera2024funsearch,novikov2025alphaevolve,assumpcao2025codeevolve}. Наша робота досліджує додаткове операційне питання: як часто ранжування оцінювача витримує незалежну перевірку та як змінюється відповідь, коли час вимірюють послідовно? Результати є обмеженим аудитом одного фреймворку й родини моделей, а не загальним бенчмарком coding agents.

% =====================================================================
\section{Обмеження та загрози валідності}\label{sec:limitations}

\begin{itemize}
\item \textbf{Обсяг задач.} Дослідження використовує дві невеликі сконструйовані задачі. Пакування кіл -- відомий бенчмарк, а не нова задача; наближений kNN є спеціально створеною задачею. Результати не встановлюють якість на реальних помилках репозиторіїв чи науковому коді.
\item \textbf{Обсяг моделей.} Шкала розмірів охоплює кілька конфігурацій Gemma 4. Порівняння з хмарною допомогою використовує один режим Argon. Якість моделі, архітектура та середовище розгортання змінюються разом; експерименти не визначають загального закону масштабування.
\item \textbf{Неповні кампанії.} У дослідженні Colab завершено 167/179 запланованих запусків, у пілоті на ноутбуці -- 12/18. Відсутні прогони зосереджені в режимі LC, обмеженому квотою. У пілоті лише три seeds на задачу й режим.
\item \textbf{Розвідувальний висновок.} Обидві реєстрації передбачали дескриптивний аналіз. Тести в цій статті додано після отримання результатів; деякі значення $p$ не скориговано, а корекція Голма стосується лише явно визначених груп порівнянь. Малі вибірки обмежують точність.
\item \textbf{Змішування збірок моделей і середовищ.} У Colab 12B B виник збій застосування diff, залежний від збірки. У порівнянні 12B на різних пристроях відрізнялися digest моделі, версія Ollama та середовище виконання. Порівняння вартості 31B проводилося за іншої паралельності.
\item \textbf{Обмеження оцінювача і holdout.} Незалежна перевірка сильніша за робочий оцінювач, але не є доказом коректності. Вона покриває конкретні обмеження та скінченні тестові входи. Пошук може використати вимоги, не охоплені жодним оцінювачем.
\item \textbf{Чутливість часу.} Fitness kNN залежить від настінного часу. Серійне вимірювання виконали для початкової програми та всіх валідних кандидатів вище за неї, але не для кожного згенерованого кандидата. Близько 1{,}900 кандидатів нижче за початкову оцінку не мають нового вимірювання; для повторно виміряних валідних програм лишається варіативність.
\item \textbf{Доступ до даних.} Попередні реєстрації, журнали відхилень, run-level аналіз, серійний перерахунок і код відкриті. Повні журнали викликів моделей та каталоги сирих прогонів зберігаються в окремих архівах і не входять до публічного Git snapshot.
\end{itemize}

% =====================================================================
\section{Доступність даних і коду}\label{sec:availability}

Попередні реєстрації, журнали відхилень, підсумки run-level розвідувального аналізу, файли серійної повторної перевірки kNN та код доступні у відкритому snapshot репозиторію на коміті \code{57a55d3d23d6c7c3dd7be3abfc1a0d9c758fc6b5}~\citep{pylypenko2026artifacts}. Зареєстровані протоколи: \code{docs/colab\_scale/PREREGISTRATION.md} і \code{docs/openevolve\_pilot/PREREGISTRATION.md}; відхилення кампаній описані у відповідних файлах \code{DEVIATIONS.md}. Повні каталоги запусків і журнали викликів моделей збережено в окремих архівах та не поширено в цій чернетці. Рисунки згенеровано з зафіксованого run-level результату \code{docs/colab\_scale/results/exploratory\_tests\_2026-11-19.json} скриптом \code{scripts/paper\_ii/make\_figures.py}. PDF можна відтворити командою \code{bash scripts/paper\_ii/build.sh}. Архівний DOI не заявляємо.

% =====================================================================
\section{Висновки}\label{sec:conclusion}

У двох обмежених дослідженнях еволюційний пошук програм під керуванням LLM знаходив корисні рішення для пакування кіл у різних конфігураціях Gemma 4, а незалежна перевірка виявляла розбіжності між робочою оцінкою та перевіреною поведінкою. Пошук на 1{,}000 викликів знайшов сильні валідні пакування, але не завжди передавав їх далі як обраний результат. Розвідувальний серійний аналіз kNN виявив вищі оцінки у восьми з дев'яти пар MoE LC--L, поряд із часткою хибних покращень 28\% серед кандидатів LC вище за початкову оцінку. Неповні кампанії, обсяг задач, чутливість часу, збій збірки та розвідувальний характер статистики обмежують ширші висновки. Майбутні дослідження мають попередньо порівняти дизайни оцінювачів, застосовувати перевірку коректності до оцінювання продуктивності, використовувати незалежні задачі та публікувати сирі журнали, перевірені результати й повні витрати.

\section*{Прозорість дослідження}
\textbf{Фінансування та конфлікт інтересів.} Зовнішнього фінансування не заявлено. Для дослідження використано обладнання та облікові записи сервісів автора. Автор не заявляє фінансового конфлікту інтересів, пов'язаного з OpenEvolve чи цитованими системами.

\textbf{Інституційні відносини.} Автор повідомляє про статус аспіранта Державного торговельно-економічного університету та співпрацю з кафедрою інженерії програмного забезпечення та кібербезпеки. Це формулювання описує заявлений автором статус і співпрацю; воно не означає, що кафедра виконувала чи фінансувала експерименти.

\textbf{Використання інструментів ШІ.} AI-асистенти допомагали в розробці скрипта аналізу, підготовці рукопису й перекладі українською. Вони не проводили модельні кампанії. Автор відповідає за перевірку аналізу, посилань, перекладу та остаточних тверджень. Ця чернетка не проходила незалежного людського рецензування.

\bibliographystyle{plainnat}
\bibliography{../references}

\appendix
\section{Попередня реєстрація та відхилення}\label{app:registration}

Пілот на ноутбуці зареєстровано до запусків під тегом \code{oe-pilot-prereg}; дослідження масштабу на Colab зареєстровано до запусків під тегом \code{colab-scale-prereg}. Обидва протоколи передбачали лише дескриптивні підсумки й не містили тестів статистичної значущості. Основні відхилення, важливі для інтерпретації: (1) дослідження Colab зупинено після 167 із 179 запусків, коли квота не дала завершити 12 LC-прогонів; (2) режим пілоту LC закрито після вичерпання квоти, завершено 12 із 18 запланованих L/B запусків; (3) збірка Colab 12B B багаторазово створювала незастосовні diff; (4) результати часу kNN повторно виміряно послідовно після того, як паралельні вимірювання тієї самої початкової програми дали оцінки від 0.02 до 1.0; (5) інференційні тести додано після завершення кампаній. Повна хронологія наведена у пов'язаних журналах відхилень.

\section{Карта доказів}\label{app:evidence}

\begin{table}[tb]
\centering\small
\caption{Основні твердження рукопису та джерела даних у репозиторії.}
\label{tab:evidence}
\begin{tabularx}{\linewidth}{@{}P{0.33\linewidth}L@{}}
\toprule
Твердження або аналіз & Джерело \\
\midrule
Статус прогонів Colab і результати масштабу & \path{docs/COLAB_SCALE_RESULTS.md}; \path{docs/colab_scale/FINAL_RUNS_2026-11-11.json} \\
Результати пілоту & \path{docs/OPENEVOLVE_PILOT_RESULTS.md}; \path{docs/openevolve_pilot/summary_2026-11-19.json} \\
Зареєстрований дизайн і відхилення & \path{docs/colab_scale/PREREGISTRATION.md}; \path{docs/colab_scale/DEVIATIONS.md}; відповідні файли пілоту \\
Серійна корекція kNN & \path{docs/colab_scale/results/knn_serial_retiming_2026-11-19.jsonl}; \path{knn_serial_above_seed_2026-11-19.jsonl} \\
Розвідувальна статистика & \path{docs/colab_scale/results/exploratory_tests_2026-11-19.json}; \path{scripts/analysis/oe_exploratory_tests.py} \\
Рисунки статті & \path{scripts/paper_ii/make_figures.py}; рисунки у \path{docs/paper_ii/figures/} \\
\bottomrule
\end{tabularx}
\end{table}

\end{document}
